% ============================================================
% CHAPITRE 1 : Introduction & Modélisation
% ============================================================
\chapterbanner{Chapitre 1 --- Introduction \& Mod\'elisation}{ch1blue}

\begin{tcolorbox}[colback=ch1blue!5,colframe=ch1blue!75!black,title={Signaux d'un syst\`eme r\'egul\'e}]
\begin{itemize}
  \item $w(t)$ : consigne (setpoint, valeur d\'esir\'ee)
  \item $y(t)$ : grandeur r\'egl\'ee (sortie mesur\'ee par capteur)
  \item $x(t)$ : grandeur r\'egl\'ee brute (non mesur\'ee)
  \item $e(t) = w(t) - y(t)$ : signal d'erreur (d\'eviation)
  \item $u(t)$ : grandeur de commande (sortie r\'egulateur)
  \item $v(t)$ : perturbation (entr\'ee non d\'esir\'ee)
  \item $n(t)$ : bruit de mesure (sur le capteur)
\end{itemize}
\end{tcolorbox}

\begin{tcolorbox}[colback=ch1blue!5,colframe=ch1blue!75!black,title={Sch\'ema bloc canonique}]
\begin{center}
\begin{tikzpicture}[>=Stealth, scale=0.52, transform shape,
  block/.style={draw, minimum height=0.72cm, minimum width=1.45cm, font=\tiny, align=center},
  yellowblock/.style={block, fill=yellow!35},
  violetblock/.style={block, fill=blue!20},
  sum/.style={draw, circle, inner sep=1pt, minimum size=0.45cm, font=\tiny},
  signal/.style={font=\tiny, align=center}]
  \node[sum] (sum) at (0,0) {$\Sigma$};
  \node[violetblock, right=0.72 of sum] (Gr) {R\'egulateur};
  \node[yellowblock, right=1.52 of Gr, minimum width=2.2cm] (Ga) {Amplificateur,\\actionneur};
  \node[yellowblock, right=0.45 of Ga] (Gp) {Processus};
  \node[yellowblock, right=0.45 of Gp] (Gm) {Capteur};
  \coordinate[right=0.95 of Gm] (out);
  \node[draw=red!75!black, rounded corners=2pt, inner xsep=4pt, inner ysep=6pt, fit=(Ga) (Gp) (Gm)] (sys) {};
  \node[signal, xshift=-20pt,above=2pt of sys,red!75!black] {Syst\`eme \`a r\'egler};

  \draw[->] (-1.7,0) node[left, signal]{Consigne\\$w(t)$} -- (sum);
  \draw[->] (sum) -- node[above=2pt, signal]{Erreur\\$e(t)$} (Gr);
  \draw[->] (Gr) -- node[pos=0.4,above=2pt, signal]{Commande\\$u(t)$} (Ga);
  \draw[->] (Ga) -- (Gp);
  \draw[->] (Gp) -- (Gm);
  \draw[->] (Gm) -- (out) node[right, signal]{Grandeur r\'egl\'ee\\mesur\'ee $y(t)$};
  \draw[->] (out) |- ([yshift=-0.1cm]Gm.south) -| (sum) node[near start, below, signal] {$-$};
  \draw[->] ($(Gp.north)+(0,1.1)$) node[above, signal]{Perturbations\\$v(t)$} -- (Gp.north);
\end{tikzpicture}
\end{center}
\vars{Le bloc \og Syst\`eme \`a r\'egler\fg{} regroupe l'amplificateur/actionneur, le processus et le capteur.}
\end{tcolorbox}

\begin{tcolorbox}[colback=ch1blue!5,colframe=ch1blue!75!black,title={Modes de r\'egulation (correspondance / maintien)}]
\begin{itemize}
  \item \textbf{Maintien} (r\'egulation) : $w = \text{cst}$, rejeter $v(t)$
  \item \textbf{Poursuite} (asservissement) : $y(t)$ suit $w(t)$ variable
  \item \textbf{Manuel} : op\'erateur contr\^ole l'actionneur
  \item \textbf{Automatique} : le r\'egulateur agit sur $e(t)$
\end{itemize}
\end{tcolorbox}

\begin{tcolorbox}[colback=ch1blue!5,colframe=ch1blue!75!black,title={Lin\'earit\'e \& Superposition (LTI)}]
\begin{itemize}
  \item \textbf{Homog\'en\'eit\'e} : $u \to y \Rightarrow \alpha u \to \alpha y$
  \item \textbf{Additivit\'e} : $u_1 \to y_1,\; u_2 \to y_2 \Rightarrow u_1{+}u_2 \to y_1{+}y_2$
  \item \textbf{Invariance temporelle} : $u(t{-}T) \to y(t{-}T)$
\end{itemize}
\vars{\textbf{Superposition} : la r\'eponse totale = somme des r\'eponses \`a chaque entr\'ee prise isol\'ement (les autres mises \`a 0).}
\end{tcolorbox}

\begin{tcolorbox}[colback=ch1blue!5,colframe=ch1blue!75!black,title={Mod\'elisation par \'eq. diff\'erentielle}]
Forme g\'en\'erale d'ordre $n$ :
\tightmath{
\sum_{k=0}^{n} a_k \frac{d^k y}{dt^k} = \sum_{k=0}^{m} b_k \frac{d^k u}{dt^k}
}
$n \geq m$ (syst\`eme causal, degr\'e relatif $d = n - m$).
\begin{center}
\renewcommand{\arraystretch}{1.05}
\scriptsize
\begin{tabular}{@{}lll@{}}
\toprule
Syst\`eme & \'Eq. diff. & $\tau$ \\
\midrule
Circuit RC & $RC\,\dot{y} + y = u$ & $RC$ \\
Circuit RL & $\frac{L}{R}\dot{i} + i = u/R$ & $L/R$ \\
R\'eservoir & $A\,\dot{h} + \frac{h}{R_h} = q_{in}$ & $A \cdot R_h$ \\
Thermique & $C_{th}\dot{T} + \frac{T}{R_{th}} = P$ & $C_{th}R_{th}$ \\
Moteur CC & $J\dot{\omega} + f\omega = K_t i$ & $J/f$ \\
Masse-amort. & $m\dot{v} + R_f v = F$ & $m/R_f$ \\
\bottomrule
\end{tabular}
\end{center}
\end{tcolorbox}

\begin{tcolorbox}[colback=ch1blue!5,colframe=ch1blue!75!black,title={Gain statique}]
\tightmath{
\boxed{K = \frac{\Delta y_\infty}{\Delta u_\infty} = \lim_{s \to 0} G(s) = \frac{b_0}{a_0}}
}
\end{tcolorbox}


\begin{tcolorbox}[colback=ch1blue!5,colframe=ch1blue!75!black,title={Sch\'ema bloc avec int\'egrateurs (\'eq. diff $\to$ bloc)}]
Pour $\ddot y + a_1 \dot y + a_0 y = b_0\,u$ (ordre 2) :
\begin{itemize}
  \item Isoler la d\'eriv\'ee la plus haute : $\ddot y = b_0 u - a_1 \dot y - a_0 y$
  \item Cha\^iner des int\'egrateurs : $\ddot y \xrightarrow{\int} \dot y \xrightarrow{\int} y$
  \item Retrouver $\ddot y$ via sommateur + r\'etroactions $(-a_0, -a_1)$
  \item Gain d'entr\'ee $= b_0$ (ou $b_m$ sur la d\'eriv\'ee d'ordre $m$ de $u$)
\end{itemize}
\begin{center}
\begin{tikzpicture}[>=Stealth, scale=0.92, transform shape,
  block/.style={draw, minimum height=0.82cm, minimum width=1.15cm, align=center},
  gain/.style={draw, minimum height=0.65cm, minimum width=0.9cm, align=center},
  sum/.style={draw, circle, inner sep=1pt, minimum size=0.48cm}]
  
  \node[sum] (sum) at (0,0) {$\Sigma$};
  \node[gain, left=0.8 of sum] (b0) {$b_0$};
  \node[block, right=0.6 of sum] (int1) {$\dfrac{1}{s}$};
  \node[block, right=0.6 of int1] (int2) {$\dfrac{1}{s}$};
  \node[coordinate, right=0.7 of int2] (out) {};

  \node[gain] (a1) at (1.35,-1.15) {$a_1$};
  \node[gain] (a0) at (2.95,-1.55) {$a_0$};

  \draw[->] (-2.3,0) node[left] {$u$} -- (b0);
  \draw[->] (b0) -- (sum);
  \draw[->] (sum) -- node[above] {$\ddot y$} (int1);
  \draw[->] (int1) -- node[above] {$\dot y$} coordinate[pos=0.5] (ydotmid) (int2);
  \draw[->] (int2) -- node[above] {$y$} coordinate[pos=0.5] (ymid) (out);

  \draw[->] (ydotmid) |- (a1);
  \draw[->] (a1) -| (sum.south);

  \draw[->] (ymid) |- (a0);
  \draw[->] (a0) -| (sum.south);

  \node[font=\tiny] at (-0.3,0.2) {$+$};
  \node[font=\tiny] at (+0.2,-0.35) {$-$};
\end{tikzpicture}

\end{center}
\astuce{Forme g\'en\'erale : nb int\'egrateurs = ordre $n$ du d\'enominateur.}
\end{tcolorbox}

\begin{tcolorbox}[colback=ch1blue!5,colframe=ch1blue!75!black,title={P\^oles, p\^ole dominant \& dur\'ee du r\'egime transitoire}]
\textbf{P\^oles} $p_i$ = racines du d\'enominateur de $G(s)$ :
\begin{itemize}
  \item R\'eels $p_i < 0$ : mode ap\'eriodique, $e^{p_i t}$
  \item Complexes $p_{1,2} = -\delta \pm j\omega_0$ : oscillatoire, $e^{-\delta t}\sin(\omega_0 t)$
  \item $\text{Re}(p_i) \geq 0$ : syst\`eme instable !
\end{itemize}
\textbf{P\^ole dominant} = p\^ole avec $|\text{Re}(p_i)|$ le plus petit (le plus lent) :
\tightmath{
\boxed{T_{\text{r\'eg}} = \frac{3}{|\text{Re}(p_{\text{dom}})|}} \qquad T_0 = \frac{2\pi}{|\text{Im}(p_{\text{dom}})|}
}

\vars{$t_r$ : temps de mont\'ee. $T_{\text{r\'eg}}$ : temps pour entrer \emph{et rester} dans la bande $\pm5\%$ ; le facteur $3$ est d\'ej\`a inclus dans la formule. $T_0$ : p\'eriode des oscillations.}
\vars{Si un p\^ole est en $s=0$, alors c'est lui le p\^ole dominant : il est plus lent que tout p\^ole de partie r\'eelle strictement n\'egative.}
\end{tcolorbox}

\begin{tcolorbox}[colback=ch1blue!5,colframe=ch1blue!75!black,title={Astuces \& Pi\`eges Ch.1}]
\piege{\textbf{Gain statique $\neq$ gain permanent} : gain stat. $= \lim_{s\to 0}G(s)$ (peut valoir $\infty$ si $\alpha{\ge}1$). Gain permanent $K{=}\lim_{s\to 0}s^\alpha G(s)$ (toujours fini).}\\[1pt]
\piege{\textbf{Ordre $n$ vs type $\alpha$} : $n$ = degr\'e du d\'enom. ; $\alpha$ = nb de p\^oles en $s{=}0$.}\\[1pt]
\piege{\textbf{P\^ole dominant} : ignorer les p\^oles $\gg 10\times$ plus rapides. Un z\'ero proche d'un p\^ole l'annule (simplification p\^ole-z\'ero).}\\[1pt]
\astuce{Pour $G(s)=\frac{7}{s(s{+}1)(s{+}10)}$ : $\alpha{=}1$, gain perm. $=7/10$, gain stat.$=\infty$, p\^ole dom. $s{=}0$, $T_{\text{r\'eg}}\to\infty$.}
\end{tcolorbox}
